\chapter*{Preface}
\addcontentsline{toc}{chapter}{Preface}

This book is about the gap between an association and an answer. A regression
coefficient is an association. The question that motivated the regression is
almost always causal: whether a programme raised earnings, whether a drug
shortened a hospital stay, whether a tax changed behaviour. Closing that gap
requires assumptions that the data cannot supply, and the discipline of causal
inference is largely the discipline of stating those assumptions precisely
enough that a reader can disagree with them.

The organising idea is a separation of concerns. Part~\ref{part:identification}
asks whether the causal quantity is a function of the observed distribution at
all. This is identification, and it is a question about assumptions rather than
about sample size. Part~\ref{part:estimation} asks how to estimate an identified
quantity efficiently and how to do so without betting everything on one
parametric model. Part~\ref{part:design} asks what to do when the identifying
assumptions are doubtful, which is the normal case.

The separation matters because the three failures look alike in a data analysis
and have nothing to do with one another. An unidentified estimand cannot be
rescued by a larger sample. An inefficient estimator of an identified quantity
is a waste rather than an error. A well estimated quantity under a false
assumption is the most dangerous of the three, because its confidence interval
is narrow and centred in the wrong place.

A longer treatment of the philosophical background is given in
\citet{achterberg2018foundations}, and readers who want the graphical language
developed from first principles should read \citet{ferreiro2007graphs}
alongside Part~\ref{part:identification}.

The prerequisites are a first course in mathematical statistics and comfort with
conditional expectation. Measure theoretic care is used where it changes an
answer and suppressed where it does not. Every theorem is proved, and the proofs
are short by design: a proof that runs for four pages usually indicates that the
statement is doing too much work.

Readers coming from applied practice may prefer to begin with
Chapter~\ref{chap:propensity} and return to Part~\ref{part:identification} when
a weighting estimator misbehaves, which it will. Readers coming from
econometrics will find Chapter~\ref{chap:design} familiar and
Chapter~\ref{chap:graphical} less so, and the connection between the two
languages is drawn explicitly in Section~\ref{sec:design-dag-view}.

The fourth edition rewrites the treatment of doubly robust estimation in
Chapter~\ref{chap:doubly-robust} around the sample splitting argument, which
makes the rate condition transparent, and adds the sensitivity contours of
Chapter~\ref{chap:sensitivity}. The exercises have been removed. They were never
good, and the space is better spent on the material.
